\centerline{\bf \Huge Перебор случаев}
\title{Перебор случаев}


\problem{1} C понедельника по среду гном ест на завтрак манную кашу, с четверга по субботу — рисовую кашу, а в воскресенье делает себе яичницу. По чётным числам месяца гном говорит правду, а по нечётным — неправду. В какие из первых десяти дней августа 2016 года он мог сказать: «Завтра я буду есть на завтрак манную кашу»? Обоснуйте ваш ответ.

\problem{2} Кузнечик прыгает вдоль прямой вперёд на 80 см или назад на 50 см. Может ли он менее чем за 7 прыжков удалиться от начальной точки ровно на 1 м 70 см?

\problem{3} Петя обратил внимание, что дата проведения Турнира Архиме- да, записанная восемью цифрами (22.01.2012) обладает интересной особенностью: переставив первые четыре цифры, можно получить номер года. А какие ещё даты в этом году имеют такое же свойство?

\problem{4} У бабушки три внука. Если внук заканчивал первый класс, то ба- бушка дарила ему одну книгу, если заканчивал второй класс, то бабушка дарила ему две книги, если третий класс, то три книги и т. д. Книги, полученные в подарок за все годы, внуки ставили на одну полку. Сейчас на полке 23 книги. Известно, что один из внуков старше остальных не меньше чем на два года. Какой класс он окончил? 

\problem{5} В коробке лежат синие, красные и зелёные карандаши — всего 20 штук. Синих карандашей в 6 раз меньше, чем зелёных. Красных карандашей тоже меньше, чем зелёных. Сколько карандашей нужно вытащить из коробки, чтобы вероятность того, что среди вытащенных карандашей есть хотя бы один красный, была равна 1?

\problem{6} На гранях кубика расставлены числа от 1 до 6. Кубик бросили два раза. В первый раз сумма чисел на четырёх боковых гранях оказалась равна 12, во второй — 15. Какое число написано на грани, противоположной той, где написана цифра 3?